% !PRINTED 618
\weberpage{618}
%% [page furniture, not body text -- source-confirmed 2026-08-10]
%% 618 \quad Dreiundzwanzigster Abschnitt. \quad \S\,168.

Die erste der Bedingungen (9) fordert, da\ss\ $p$ durch die Hauptform der Diskriminante $D$ eigentlich darstellbar sei (\S\,122), die zweite, da\ss
\weq{10}{p \equiv 1 \pmod{m}}
sei.

Damit ein Primideal $\mathfrak{p}$ ersten Grades in $\Omega$, das nicht in $S$ aufgeht, in $\bar{E}A$ enthalten sei, ist notwendig und hinreichend, da\ss

1. $\mathfrak{p}$ ein Hauptideal, d.\,h. eine existierende Zahl $\pi$ sei, und da\ss\ diese Zahl $\pi$ in $A$ enthalten sei, d.\,h. da\ss
\[
\text{2. } N(\pi) = p \equiv 1 \pmod{m}
\]
sei. Diese Bedingungen stimmen genau mit den Bedingungen (9) \"uberein, und folglich ist $\mathfrak{K}(D, \varrho)$ \gesperrt{der Klassenk\"orper der Gruppe $A$.}

Der Grad des K\"orpers $\mathfrak{K}(A)$ ist daher nicht nur h\"ochstens, sondern genau gleich $h$.

Die Zahl $(Z, R_m)$ ist $=1$, wenn in $m$ keine der charakteristischen Primzahlen des \S\,108 aufgeht, und dann ist der Grad des K\"orpers $\mathfrak{K}(A)$ gleich $H'\varphi(m)$. Nehmen wir umgekehrt alle charakteristischen Primzahlen und Primzahlpotenzen in $m$ auf, so wird $R_m = R$ \gesperrt{die Gruppe der absoluten Normenreste}, und es ist
\[
(Z, R) = 2g,
\]
wenn $g$ die Anzahl der Geschlechter der Diskriminante $D$ ist (\S\,113, 7.). Bezeichnet also $H'_g$ die Anzahl der Klassen eines Geschlechtes, so ist $H' = g H'_g$, und die Formel (7) ergibt:
\weq{11}{h = \tfrac{1}{2} H'_g \varphi(m).}

Nach \S\,138 zerf\"allt die Klassenfunktion $H_{-D}(u)$ durch Adjunktion von Quadratwurzeln in so viele Faktoren, als es Geschlechter von Formenklassen gibt, und der Grad eines jeden dieser Faktoren ist $H'_g$.

W\"are nun die Klassengleichung durch Adjunktion irgend welcher Einheitswurzeln noch weiter zerlegbar, als nach den Geschlechtern, so k\"onnte man $m$ so annehmen, da\ss\ auch diese Einheitswurzeln und auch $\sqrt{\Delta}$ in dem K\"orper $\mathfrak{R}(\varrho)$ enthalten w\"aren.

Ist dann $h''$ der niedrigste Grad, auf den die Klassengleichung durch Adjunktion von Einheitswurzeln reduziert werden kann, so ist der absolute Grad des K\"orpers $\mathfrak{K}(D, \varrho)$ h\"ochstens gleich $h''\varphi(m)$
%CONT>>
